\documentclass[11pt]{article}
\usepackage[margin=1in]{geometry}
\usepackage{amsmath,amssymb}
\usepackage{natbib}
\usepackage{graphicx}
\usepackage{booktabs}
\usepackage{placeins}
\usepackage[font=small,labelfont=bf]{caption}
\usepackage[colorlinks=true,allcolors=blue]{hyperref}
\setlength{\emergencystretch}{3em}
\newcommand{\rr}{\langle r\rangle}

\title{\bfseries Spurious level repulsion in geophysical random-matrix analysis:\\
declustering, seasonal gating, and the Wishart null}
\author{Ruqing Chen\\[2pt]
\small GUT Geoservice Inc., Montr\'eal, Qu\'ebec, Canada\\
\small \texttt{ruqing.chen@gutgeoservice.com}}
\date{\today}

\begin{document}
\maketitle

\begin{abstract}
\noindent
The consecutive spacing ratio $\rr$ is increasingly used to claim Wigner--Dyson
(GOE) ``level repulsion'', and hence organisation, in geophysical and climate
data. But $\rr$ is two-ended, and several mundane data structures mimic GOE
without any underlying repulsion. We document three such mechanisms with
controlled demonstrations and give the correct controls. (1)~\emph{Declustering:}
removing points closer than a fraction of the median spacing drives even pure
Poisson points to GOE-like $\rr$, manufacturing regularity from randomness.
(2)~\emph{Seasonal gating:} when events can occur only in a fixed season, the
inter-event interval distribution is bimodal and the naive $\rr$ lands between
Poisson and GOE as an artifact of that marginal. For 41 major sudden
stratospheric warmings the naive interval $\rr=0.48$ is permutation-invariant
(scalar), while the correct de-seasonalised test --- the per-winter event count
--- shows mild \emph{under-dispersion} (variance/mean $=0.67$, 95\% CI excluding
unity), i.e.\ a weak refractory clock, not level repulsion. (3)~\emph{The Wishart
null:} the eigenvalue spacing of a sample correlation matrix is GOE for
independent data, so ``GOE level spacing'' in a climate correlation matrix is the
finite-sampling null; the physical signal lies in the eigenvalue \emph{density}
(departures from Marchenko--Pastur), not in its spacing. The honest diagnostics
are raw $\rr$ with the coefficient of variation, the permutation (shuffle) test,
an appropriate physical null, and --- for correlation matrices --- the eigenvalue
density rather than its spacing.
\end{abstract}

\section{Introduction}

Random-matrix theory (RMT) provides a celebrated test for ``order'' in a
sequence of levels: the spacings between consecutive levels follow Poisson
statistics for an uncorrelated sequence and the Wigner--Dyson statistics of the
Gaussian ensembles for a spectrum with level repulsion \citep{mehta2004}. A
convenient, unfolding-free measure is the mean ratio of consecutive spacings
\citep{oganesyan2007,atas2013},
\begin{equation}
\rr = \Big\langle \frac{\min(s_i,s_{i+1})}{\max(s_i,s_{i+1})}\Big\rangle,
\qquad s_i = x_{i+1}-x_i,
\label{eq:r}
\end{equation}
which takes the reference values $0.386$ (Poisson), $0.531$ (GOE) and $0.603$
(GUE). The statistic is attractive precisely because it needs no unfolding, and
it is now applied well beyond quantum spectra --- to earthquake catalogues,
eruption sequences, climate indices and ore-deposit records --- often to argue
that a geophysical system shows ``level repulsion'' and is therefore organised by
a single coherent process.

This note is a caution. The ratio $\rr$ is two-ended: a value above the Poisson
reference can arise from genuine level repulsion, but equally from a narrow or
structured marginal spacing distribution with no sequential repulsion at all, or
from pre-processing. We document three distinct ways in which geophysical and
climate analyses obtain GOE-like $\rr$ without any underlying repulsion, each
with a controlled demonstration, and we state the diagnostic that exposes it. The
first mechanism was reported in a companion study of orogenic-gold vein spacing
\citep{chen2026sigma}; we recapitulate it and add two more that are specific to
climate data, unifying all three under a single rule: \emph{compare the observed
statistic against the permutation of its own data and against the correct null,
never against the textbook GOE value alone.}

\section{The ratio statistic and the permutation test}

Genuine level repulsion has two fingerprints: a marginal spacing distribution
that is depleted at small spacings (so that $P(s\!\to\!0)\to0$), \emph{and}
sequential anti-correlation (spectral rigidity). The ratio $\rr$ responds to
both. The decisive control is the \emph{permutation (shuffle) test}: randomly
re-order the spacings, which preserves the marginal distribution exactly but
destroys any sequential correlation, and recompute $\rr$. If
$\rr_{\mathrm{obs}}\approx\rr_{\mathrm{shuf}}$, the value is a property of the
marginal alone --- it carries no sequential repulsion, however far above Poisson
it lies \citep{chen2026nucleosome}. We also report the coefficient of variation
$\mathrm{CV}=\sigma_s/\bar s$, which is $1$ for Poisson, $<1$ for regular sets and
$>1$ for clustered sets. With these two numbers, and an appropriate null, each of
the three hazards below is immediately diagnosable.

\section{Hazard 1: declustering fabricates GOE}

A common pre-processing step is \emph{declustering}: removing points that lie
closer together than some fraction of the median spacing. This manufactures
regularity from any input. Figure~\ref{fig:hazards}a shows the effect on pure
Poisson points: as the minimum-separation threshold increases from $0$ to $0.7$
median spacings, $\rr$ rises monotonically from the Poisson value $0.386$ through
the GOE value $0.531$ to $0.57$. No repulsion is present in the data; the GOE-like
value is entirely an artifact of the thinning rule. Mechanistically,
declustering injects an \emph{artificial hard core}: by deleting the smallest
spacings it forces $P(s\!\to\!0)\to0$, the very depletion that defines level
repulsion, so it manufactures a steric-exclusion signature the data do not
contain --- the inverse of the genuine hard cores documented in nucleosome
arrays and ice-wedge networks. This hazard was identified,
with a field example, in a study of orogenic-gold vein spacing
\citep{chen2026sigma}, where raw vein positions are clustered ($\rr=0.34$) but
rise to $0.44$--$0.50$ under declustering in lockstep with random controls. The
remedy is simple: never decluster before computing a spacing statistic; report
the raw $\rr$ with CV.

\section{Hazard 2: seasonal gating}

Many geophysical events can occur only within a fixed window of the year ---
winter storms, monsoon bursts, stratospheric warmings. This seasonal gating
imprints a strong periodic structure on the sequence of event times that has
nothing to do with the system's internal dynamics, yet contaminates the naive
interval statistic.

We illustrate with major sudden stratospheric warmings (SSWs), the abrupt
reversals of the winter polar vortex. We use the 41 central dates of major
Northern-Hemisphere SSWs in the JRA-55 reanalysis (1958--2023) from the NOAA SSW
Compendium \citep{butler2017,charlton2007}. Computed naively on the calendar
sequence, the inter-event interval ratio is $\rr=0.484$ --- comfortably between
Poisson and GOE, and easily mis-read as partial ``repulsion''. It is nothing of
the kind. The interval distribution is strongly bimodal
(Fig.~\ref{fig:hazards}b): $85\%$ of intervals exceed $270$ days (the gaps across
each SSW-free summer) and only a handful fall within a winter. The permutation
test confirms the artifact: $\rr_{\mathrm{shuf}}=0.475$, indistinguishable from
the observed value, so the elevation above Poisson is a property of the bimodal,
seasonally-imposed marginal --- a scalar, not sequential repulsion. The
intermediate value follows directly from the bimodality: consecutive
cross-summer gaps (the dominant mode) give ratios near unity, while the
occasional pairing of a within-winter interval with a cross-summer gap gives
ratios near zero, and the average of these extremes falls between the Poisson and
GOE references. Because this is a property of the marginal alone, the permutation
test reproduces it exactly.

The physically meaningful question --- does the vortex resist a second warming
(a refractory period) or fire at random? --- is answered only after
de-seasonalising. Counting major SSWs per winter season over the 66 winters
spanned (Fig.~\ref{fig:ssw}) gives a mean of $0.62$ and a variance of $0.42$, a
variance-to-mean ratio of $0.67$ with a bootstrap $95\%$ confidence interval of
$[0.49,0.86]$ that excludes unity. The counts are therefore mildly
\emph{under-dispersed}: winters with two SSWs are rarer, and winters with exactly
one are commoner, than a Poisson process would produce, consistent with a
recovery (refractory) timescale after each event \citep{ineson2024}. This is a
weak scalar refractory clock --- the opposite end of the spacing axis from
clustering, but not Wigner--Dyson level repulsion, and invisible to the naive
interval $\rr$.

\section{Hazard 3: the Wishart null}

RMT is also applied to climate \emph{correlation matrices}: the $P\times P$
matrix of correlations among $P$ spatial cells or indices, whose eigenvalues are
then examined for GOE level statistics, taken as evidence of collective
organisation \citep{santhanam2001}. Here the trap is more basic. The eigenvalue
spacing of a sample correlation (Wishart) matrix is GOE \emph{by construction},
independent of any real correlation in the data. Figure~\ref{fig:hazards}c
demonstrates this: for $P$ \emph{independent} Gaussian time series of length $T$
--- zero true correlation --- the bulk eigenvalue ratio of the sample correlation
matrix is $\rr\approx0.53$ across a range of $(P,T)$, squarely at the GOE value.
GOE-class level spacing is thus the finite-sampling \emph{null} for any
correlation matrix; reporting it as a climate finding is to report the null. The
bulk spacing is fixed by random-matrix universality --- the local eigenvalue
statistics shared by essentially every large matrix in the ensemble, independent
of the entries' physical origin --- so it tests that universality class, not any
climate or geophysical dynamics.

The genuine RMT signal in a correlation matrix is not in the spacing but in the
eigenvalue \emph{density}: eigenvalues that escape the Marchenko--Pastur support
expected for pure noise mark real collective modes \citep{marchenko1967}, exactly
as established for financial covariance matrices \citep{laloux1999,plerou2002}.
The level spacing of the bulk is GOE whether or not those modes exist. The remedy
is to test the eigenvalue density against Marchenko--Pastur, and reserve any
spacing analysis for the de-noised spectrum.

\begin{table}[t]
\centering
\caption{Three routes to spurious GOE-like $\rr$ and the diagnostic that exposes
each.}
\label{tab:hazards}
\begin{tabular}{lll}
\toprule
Hazard & Spurious result & Correct diagnostic\\
\midrule
Declustering        & random points $\to\rr\approx$ GOE      & raw $\rr$, no thinning; CV\\
Seasonal gating     & bimodal intervals $\to\rr=0.48$        & permutation test; de-seasonalised counts\\
Wishart sampling    & iid data $\to$ bulk $\rr\approx$ GOE   & eigenvalue density vs Marchenko--Pastur\\
\bottomrule
\end{tabular}
\end{table}

\begin{figure}[t]
\centering
\includegraphics[width=\textwidth]{fig_hazards.pdf}
\caption{Three ways to fabricate GOE-like level repulsion.
\textbf{(a)} Minimum-separation declustering drives pure Poisson points from
$\rr=0.386$ through the GOE value to $0.57$.
\textbf{(b)} Sudden-stratospheric-warming inter-event intervals are bimodal
(note the cross-summer gaps near and beyond one year); the naive $\rr=0.48$
equals its permutation value and is a seasonal artifact.
\textbf{(c)} The bulk eigenvalue ratio of a sample correlation matrix built from
\emph{independent} Gaussian series sits at the GOE value for every $(P,T)$ tested:
GOE spacing is the Wishart null.}
\label{fig:hazards}
\end{figure}

\begin{figure}[t]
\centering
\includegraphics[width=0.6\textwidth]{fig_ssw_counts.pdf}
\caption{De-seasonalised sudden stratospheric warmings: the per-winter event
count is mildly under-dispersed relative to Poisson (variance/mean $=0.67$, 95\%
CI $[0.49,0.86]$, excluding unity) --- a weak refractory clock, not level
repulsion. Winters with exactly one SSW are commoner, and winters with two rarer,
than Poisson predicts.}
\label{fig:ssw}
\end{figure}

\FloatBarrier
\section{Correct practice}

The three hazards share a single remedy, summarised in Table~\ref{tab:hazards}.
Before interpreting any $\rr$ above the Poisson value as level repulsion: (i)~use
the raw data, applying no declustering or thinning; (ii)~run the permutation
(shuffle) test --- if $\rr_{\mathrm{obs}}\approx\rr_{\mathrm{shuf}}$ the value is a
marginal (scalar) property, not sequential repulsion; (iii)~report the
coefficient of variation alongside $\rr$, since clustered sets ($\mathrm{CV}>1$)
can sit near or below Poisson while a narrow marginal ($\mathrm{CV}<1$) lifts
$\rr$ without repulsion; (iv)~for periodically gated events, de-seasonalise (e.g.\
analyse per-season counts) before drawing dynamical conclusions; and (v)~for
correlation matrices, test the eigenvalue density against Marchenko--Pastur
rather than reading GOE off the bulk spacing. Genuine Wigner--Dyson repulsion
survives all of these; the spurious cases do not.

\section{Conclusion}

GOE-like values of the consecutive spacing ratio are easy to obtain by accident.
Declustering manufactures them from random points; seasonal gating manufactures
them from the calendar; and the eigenvalue spacing of any sample correlation
matrix is GOE for independent data. None of these is level repulsion. The
permutation test, the coefficient of variation, de-seasonalisation, and the
Marchenko--Pastur density test distinguish a real Wigner--Dyson spectrum from its
common impostors. Applied to sudden stratospheric warmings, the honest analysis
finds not level repulsion but a weak refractory clock --- a small, real, and
correctly-attributed result.

\section*{Data and code availability}
The SSW central dates are from the NOAA SSW Compendium \citep{butler2017}
(\url{https://csl.noaa.gov/groups/csl8/sswcompendium/}). All demonstrations are
synthetic or use these public dates; code and derived tables are at
\url{https://github.com/Ruqing1963/spurious-level-repulsion} and archived on
Zenodo (DOI assigned on release).

\begin{thebibliography}{99}

\bibitem[Atas et al.(2013)]{atas2013}
Y.~Y. Atas, E.~Bogomolny, O.~Giraud, and G.~Roux.
Distribution of the ratio of consecutive level spacings in random matrix
ensembles. \textit{Phys. Rev. Lett.} \textbf{110}, 084101 (2013).

\bibitem[Butler et al.(2017)]{butler2017}
A.~H. Butler, J.~P. Sjoberg, D.~J. Seidel, and K.~H. Rosenlof.
A sudden stratospheric warming compendium.
\textit{Earth Syst. Sci. Data} \textbf{9}, 63--76 (2017).

\bibitem[Charlton \& Polvani(2007)]{charlton2007}
A.~J. Charlton and L.~M. Polvani.
A new look at stratospheric sudden warmings. Part I: Climatology and modeling
benchmarks. \textit{J. Climate} \textbf{20}, 449--469 (2007).

\bibitem[Chen(2026a)]{chen2026sigma}
R.~Chen.
Time--space asymmetry in orogenic gold systems: temporal level repulsion but
spatial clustering of ore-stage veins at the Sigma deposit. Zenodo,
\url{https://doi.org/10.5281/zenodo.20778697} (2026).

\bibitem[Chen(2026b)]{chen2026nucleosome}
R.~Chen.
The consecutive spacing-ratio of nucleosome arrays: high regularity from steric
renewal, not level repulsion. Zenodo, \url{https://doi.org/10.5281/zenodo.20839558}
(2026).

\bibitem[Ineson et al.(2024)]{ineson2024}
S.~Ineson, et al.
Statistics of sudden stratospheric warmings using a large model ensemble.
\textit{Atmos. Sci. Lett.} \textbf{25}, e1202 (2024).

\bibitem[Laloux et al.(1999)]{laloux1999}
L.~Laloux, P.~Cizeau, J.-P. Bouchaud, and M.~Potters.
Noise dressing of financial correlation matrices.
\textit{Phys. Rev. Lett.} \textbf{83}, 1467--1470 (1999).

\bibitem[Marchenko \& Pastur(1967)]{marchenko1967}
V.~A. Marchenko and L.~A. Pastur.
Distribution of eigenvalues for some sets of random matrices.
\textit{Mat. Sb.} \textbf{72}, 507--536 (1967).

\bibitem[Mehta(2004)]{mehta2004}
M.~L. Mehta. \textit{Random Matrices}, 3rd ed. Elsevier, Amsterdam (2004).

\bibitem[Oganesyan \& Huse(2007)]{oganesyan2007}
V.~Oganesyan and D.~A. Huse.
Localization of interacting fermions at high temperature.
\textit{Phys. Rev. B} \textbf{75}, 155111 (2007).

\bibitem[Plerou et al.(2002)]{plerou2002}
V.~Plerou, P.~Gopikrishnan, B.~Rosenow, L.~A.~N. Amaral, T.~Guhr, and
H.~E. Stanley.
Random matrix approach to cross correlations in financial data.
\textit{Phys. Rev. E} \textbf{65}, 066126 (2002).

\bibitem[Santhanam \& Patra(2001)]{santhanam2001}
M.~S. Santhanam and P.~K. Patra.
Statistics of atmospheric correlations.
\textit{Phys. Rev. E} \textbf{64}, 016102 (2001).

\end{thebibliography}

\end{document}
